\documentclass[12pt, a4paper]{report}

\usepackage[utf8]{inputenc}
\usepackage[T1]{fontenc}
\usepackage[french]{babel}
\usepackage{geometry}
\usepackage{graphicx}
\usepackage{times}
\usepackage{helvet}
\usepackage{color}
\usepackage{xcolor}
\usepackage{titlesec}
\usepackage{fancyhdr}
\usepackage{setspace}
\usepackage{hyperref}
\usepackage{tocloft}
\usepackage{caption}
\usepackage{booktabs}
\usepackage{array}
\usepackage{tabularx}
\usepackage{ragged2e}
\usepackage{microtype}
\usepackage{enumitem}
\usepackage{float}
\usepackage{pgfgantt}
\usepackage{listings}
\usepackage{tikz}
\usetikzlibrary{shapes.geometric, arrows.meta, positioning}
\usepackage{pgfplots}
\pgfplotsset{compat=1.18}
\usepackage{csquotes}

\geometry{top=2.5cm, bottom=2.5cm, left=2.5cm, right=2.5cm}

\definecolor{iaaiblue}{RGB}{0, 83, 155}
\definecolor{codegreen}{RGB}{0,128,0}
\definecolor{codegray}{RGB}{128,128,128}
\definecolor{codepurple}{RGB}{128,0,128}
\definecolor{backcolour}{RGB}{245,245,245}

\renewcommand{\familydefault}{\rmdefault}

\pagestyle{fancy}
\fancyhf{}
\fancyhead[L]{\small\textit{Salma Isser}}
\fancyhead[C]{\small\textit{Rapport de Stage}}
\fancyhead[R]{\small\textit{IAAI Academy}}
\fancyfoot[C]{\thepage}
\renewcommand{\headrulewidth}{0.4pt}
\usepackage{etoolbox}
\patchcmd{\chapter}{\thispagestyle{plain}}{\thispagestyle{fancy}}{}{}

\titleformat{\chapter}[block]{\sffamily\large\bfseries\color{black}}{\Roman{chapter}.}{0.5em}{}[\titlerule]
\titleformat{\section}[block]{\sffamily\normalsize\bfseries}{\thesection}{0.5em}{}
\titleformat{\subsection}[block]{\sffamily\normalsize\bfseries}{}{0em}{}
\titlespacing*{\chapter}{0pt}{10pt}{10pt}
\titlespacing*{\section}{0pt}{8pt}{4pt}
\titlespacing*{\subsection}{0pt}{6pt}{2pt}

\hypersetup{colorlinks=true, linkcolor=iaaiblue, urlcolor=iaaiblue, citecolor=iaaiblue}

\lstdefinestyle{mystyle}{
    backgroundcolor=\color{backcolour},
    commentstyle=\color{codegreen},
    keywordstyle=\color{magenta},
    numberstyle=\tiny\color{codegray},
    stringstyle=\color{codepurple},
    basicstyle=\ttfamily\footnotesize,
    breakatwhitespace=false,
    breaklines=true,
    captionpos=b,
    keepspaces=true,
    numbers=left,
    numbersep=5pt,
    showspaces=false,
    showstringspaces=false,
    showtabs=false,
    tabsize=2
}
\lstset{style=mystyle}

\begin{document}

\begin{titlepage}
  \thispagestyle{empty}
  \centering
  
  \vspace{1cm}
  
  \begin{minipage}[t]{0.4\textwidth}
    \flushleft
    \includegraphics[height=2cm]{images/logo_mundiapolis.png}
  \end{minipage}
  \hfill
  \begin{minipage}[t]{0.45\textwidth}
    \flushright
    \includegraphics[height=2.5cm]{images/logo_entreprise.png}
  \end{minipage}
  
  \vspace{2cm}
  
  {\sffamily\Large\bfseries International Academy of Artificial Intelligence}
  
  \vspace{1.5cm}
  
  {\sffamily\Large\bfseries RAPPORT DE STAGE}
  
  \vspace{1cm}
  
  {\normalsize Stagiaire :}
  
  \vspace{0.5cm}
  
  {\large\bfseries Mlle Salma ISSER}
  
  \vspace{1cm}
  
  {\large\bfseries Filière Informatique Appliquée}
  
  \vspace{0.5cm}
  
  {\large\bfseries\textcolor{iaaiblue}{Option Développement logiciel}}
  
  \vspace{1.5cm}
  
  {\large\bfseries Sujet du stage}
  
  \vspace{0.3cm}
  
  \noindent\rule{\textwidth}{0.6pt}
  
  \vspace{0.5cm}
  
  {\fontsize{16}{22}\selectfont\textcolor{iaaiblue}{
Étude, conception et développement d'une plateforme e-learning \\[0.3cm]
pour l'apprentissage de l'Intelligence Artificielle~: IAAI e-learning 101}}
  
  \vspace{0.5cm}
  
  \noindent\rule{\textwidth}{0.6pt}
  
  \vspace{1.5cm}
  
  \begin{flushleft}
    {\normalsize Encadrement :}
    
    \vspace{0.5cm}
    
    \begin{tabular}{ll}
      \textbf{M. Naim {\hspace{4cm}}} & \textbf{: Encadrant Entreprise} \\[0.5cm]
      \textbf{Mme Ichraq ESSADEQ\hspace{1cm}} & \textbf{: Encadrante Académique} \\
    \end{tabular}
    
    \vspace{1cm}
    
    {\normalsize Durée du stage : 08 mai 2026 -- 08 juillet 2026}
    
    \vspace{0.5cm}
    
    {\normalsize Lieu du stage : IAAI Academy, Casablanca}
  \end{flushleft}
  
  \vfill
  
  {\large\bfseries\textcolor{iaaiblue}{Année Universitaire : 2025-2026}}
  
  \vspace{1cm}
  
\end{titlepage}

\clearpage
\thispagestyle{empty}
\setcounter{page}{2}
\begin{center}
  {\sffamily\large\bfseries Remerciements}
\end{center}
\vspace{1cm}
\onehalfspacing
{\fontfamily{ptm}\selectfont\fontsize{12}{18}\selectfont
Au terme de ce stage au sein de l'International Academy of Artificial Intelligence (IAAI Academy), je tiens à exprimer ma profonde gratitude envers toutes les personnes qui ont contribué à la réussite de cette expérience professionnelle et à l'aboutissement du projet \textit{IAAI eLearning 101}.

Je tiens tout d'abord à remercier sincèrement \textbf{M.Naim}, mon encadrant entreprise, pour sa confiance, ses conseils précieux et son accompagnement tout au long de ce stage. Sa disponibilité et son expertise technique ont été déterminantes dans la conduite de ce projet.

Je souhaite exprimer ma reconnaissance à \textbf{Mme Ichraq ESSADEQ}, mon encadrante académique, pour son suivi rigoureux et ses conseils avisés qui m'ont permis de structurer mon travail et de progresser tout au long de ce projet.

Ma gratitude s'étend également à l'ensemble du corps enseignant de l'Université Mundiapolis dont l'enseignement de qualité m'a permis d'acquérir les compétences techniques et méthodologiques nécessaires à la réalisation de ce projet.

Enfin, je remercie ma famille et mes proches pour leur soutien indéfectible et leurs encouragements constants tout au long de mon parcours académique et professionnel.
}

\clearpage
\begin{center}
  {\sffamily\large\bfseries Résumé}
\end{center}
\vspace{0.5cm}
\onehalfspacing
{\fontfamily{ptm}\selectfont\fontsize{12}{18}\selectfont
Ce rapport présente le travail réalisé dans le cadre du stage de fin d'études effectué au sein de l'International Academy of Artificial Intelligence (IAAI Academy). Le projet porte sur la conception et le développement d'une plateforme e-learning interactive dédiée à l'apprentissage de l'Intelligence Artificielle, intitulée \textit{IAAI eLearning 101}. L'objectif principal est de démocratiser l'accès aux connaissances fondamentales en Intelligence Artificielle à travers une solution pédagogique moderne, accessible aussi bien sur le web que sur les appareils mobiles.

La plateforme s'adresse aux étudiants, aux professionnels en reconversion ainsi qu'à toute personne souhaitant découvrir les concepts de l'Intelligence Artificielle de manière progressive et structurée. Elle propose 8 modules selon un modèle freemium (premier module gratuit, modules suivants accessibles via paiement sécurisé Stripe), répartis en 57 leçons et 8 quiz. Afin de faciliter la compréhension de notions souvent abstraites, le projet s'appuie sur une approche pédagogique immersive intégrant des animations produites avec la bibliothèque Manim et narrées automatiquement. Les contenus sont proposés en français et en arabe (avec support RTL complet) afin de répondre aux besoins du public marocain et africain francophone.

La réalisation du projet a suivi une démarche structurée comprenant l'analyse des besoins, la modélisation de l'architecture fonctionnelle et de la base de données, la conception de l'interface utilisateur, ainsi que le développement de la solution à l'aide de technologies modernes telles que React, Vite, Tailwind CSS et Supabase. La plateforme intègre également un chatbot pédagogique, ARIA, destiné à accompagner l'apprenant dans son parcours. Une attention particulière a été portée à la sécurité de la plateforme : politiques RLS strictes, en-têtes CSP/HSTS, journal d'audit immuable et verrouillage anti-brute-force sur l'authentification. Le résultat obtenu est une plateforme interactive offrant une expérience d'apprentissage fluide, intuitive et sécurisée, actuellement en phase d'amélioration continue suite à la soutenance.
}

\noindent{\sffamily\normalsize\bfseries Mots-clés~:}\\[0.2cm]
{\fontfamily{ptm}\selectfont\fontsize{12}{16}\selectfont
E-learning, Intelligence Artificielle, Plateforme Interactive, React, Supabase,
Sécurité applicative, Chatbot pédagogique, Apprentissage Numérique, Animations Pédagogiques.
}

\clearpage
\begin{center}
  {\sffamily\large\bfseries Abstract}
\end{center}
\vspace{0.5cm}
\onehalfspacing
{\fontfamily{ptm}\selectfont\fontsize{12}{18}\selectfont
This report presents the work carried out as part of an end-of-studies internship at the International Academy of Artificial Intelligence (IAAI Academy). The project focuses on the design and development of an interactive e-learning platform dedicated to Artificial Intelligence education, entitled \textit{IAAI eLearning 101}. The platform provides learners with a structured, accessible, and secure educational environment, enabling them to follow organized courses across 8 modules (freemium model), complete assessments, monitor their learning progress, and obtain certificates upon successful completion of training modules.

The system was developed using a modern technology stack, including React, Vite, Tailwind CSS, Supabase, and PostgreSQL, ensuring performance, scalability, and security. Beyond the core learning features (course management, quiz evaluation with lightweight proctoring, progress tracking, certificate generation), significant effort was dedicated to hardening the platform's security posture through strict Row Level Security policies, CSP/HSTS headers, an immutable audit log, and brute-force login protection. The platform also integrates ARIA, an AI pedagogical chatbot designed to support learners throughout their courses.

This report describes the different phases of the project, from requirement analysis and system design through implementation, security hardening, and functional testing, and outlines perspectives for future improvement.
}

\noindent{\sffamily\normalsize\bfseries Keywords:}\\[0.2cm]
{\fontfamily{ptm}\selectfont\fontsize{12}{16}\selectfont
E-learning, Artificial Intelligence, Web Development, React, Supabase, PostgreSQL,
Application Security, Pedagogical Chatbot, Educational Platform.
}

\clearpage
\addcontentsline{toc}{chapter}{Liste des figures}
\listoffigures
\clearpage
\addcontentsline{toc}{chapter}{Liste des tableaux}
\listoftables
\clearpage
\addcontentsline{toc}{chapter}{Table des matières}
\tableofcontents

\clearpage
\renewcommand{\chaptername}{}
\renewcommand{\thechapter}{\Roman{chapter}}

\chapter{Introduction}
\onehalfspacing
\section{Contexte général}
L'évolution rapide des technologies numériques a profondément transformé les pratiques d'apprentissage et les modes d'accès au savoir. Parmi ces transformations, les plateformes d'enseignement en ligne occupent désormais une place centrale dans les domaines de l'éducation formelle et de la formation professionnelle continue.

Parallèlement à cette révolution pédagogique, l'Intelligence Artificielle connaît un développement exponentiel et s'impose comme l'un des domaines les plus stratégiques de l'informatique contemporaine. Malgré l'engouement croissant pour cette discipline, l'accès à des ressources pédagogiques structurées, rigoureuses et adaptées aux débutants demeure un obstacle pour de nombreux apprenants, en particulier dans les pays francophones d'Afrique et du Moyen-Orient.

C'est dans ce contexte que s'inscrit le projet \textit{IAAI eLearning 101}, une plateforme e-learning dédiée à l'apprentissage des fondamentaux de l'Intelligence Artificielle, développée au sein de l'IAAI Academy.

\section{Problématique}
La démocratisation de l'Intelligence Artificielle se heurte à plusieurs verrous pédagogiques et techniques. D'une part, les concepts fondamentaux de l'IA sont souvent perçus comme abstraits et difficiles d'accès pour les débutants. D'autre part, les plateformes e-learning existantes présentent des lacunes en termes d'interactivité, de personnalisation et d'accompagnement intelligent des apprenants.

\section{Objectifs du projet}
L'objectif principal du projet est de concevoir, développer et déployer une plateforme e-learning nommée \textit{IAAI eLearning 101}. Les objectifs spécifiques sont :
\begin{itemize}
    \item Offrir un parcours pédagogique structuré et progressif pour les débutants en IA.
    \item Proposer une expérience utilisateur moderne, intuitive et responsive.
    \item Intégrer des fonctionnalités de suivi de progression et d'évaluation, avec un mécanisme léger de surveillance (proctoring) des quiz.
    \item Fournir un chatbot pédagogique intelligent (ARIA) basé sur l'API Gemini.
    \item Mettre en place un modèle économique freemium avec paiement sécurisé.
    \item Assurer la sécurité, la performance et l'évolutivité de la plateforme.
\end{itemize}

\section{Méthodologie adoptée}
Le projet a été mené selon plusieurs étapes successives :
\begin{itemize}
    \item \textbf{Analyse des besoins :} Identification des acteurs, recueil des besoins fonctionnels et non fonctionnels.
    \item \textbf{Conception :} Élaboration des diagrammes MERISE et UML, conception de la base de données.
    \item \textbf{Développement :} Implémentation avec React, Supabase, PostgreSQL, Stripe et Gemini.
    \item \textbf{Tests et validation :} Tests fonctionnels, de sécurité, de performance et de compatibilité.
    \item \textbf{Déploiement :} Mise en production sur Vercel et Supabase.
    \item \textbf{Amélioration continue :} Durcissement sécurité et corrections post-soutenance.
\end{itemize}

\section{Organisation du rapport}
Le présent rapport est structuré en cinq chapitres :
\begin{itemize}
    \item Le \textbf{Chapitre II} présente l'organisme d'accueil et la planification du stage.
    \item Le \textbf{Chapitre III} est consacré à l'analyse des besoins et à la conception.
    \item Le \textbf{Chapitre IV} détaille la phase de réalisation et les choix technologiques.
    \item Le \textbf{Chapitre V} présente les tests et la validation.
    \item La \textbf{Conclusion générale} dresse le bilan et propose des perspectives.
\end{itemize}

\chapter{Environnement de stage}
\section{Introduction}
Ce chapitre présente le cadre du stage, l'organisme d'accueil, le cahier des charges et la planification du projet.

\section{Présentation de l'entreprise et du stage}
\subsection{La présentation de l'entreprise}
L'International Academy of Artificial Intelligence (IAAI Academy) est une organisation spécialisée dans la formation, le conseil et l'accompagnement académique dans les domaines de l'IA, de la Data Science et des technologies numériques avancées. Fondée à Casablanca, elle s'inscrit dans une dynamique de développement des compétences numériques au Maroc.

Les activités de l'IAAI Academy sont organisées autour de trois pôles principaux :
\begin{itemize}
    \item \textbf{Pôle Formations :} Programmes en IA générative, Data Science, ML, Cloud Computing et MLOps.
    \item \textbf{Pôle Services :} Accompagnement des entreprises dans leur transformation numérique.
    \item \textbf{Pôle Académique :} Encadrement des étudiants, ingénieurs et chercheurs.
\end{itemize}

\subsection{La présentation du stage}
Le stage s'est déroulé dans le cadre du projet \textit{IAAI eLearning 101}, initié pour offrir une formation d'introduction à l'IA, accessible en ligne, en français et en arabe.

Le cahier des charges définissait les spécifications suivantes :
\begin{itemize}
    \item Développer une plateforme e-learning web moderne et responsive.
    \item Mettre en place un système d'authentification et de gestion de profils.
    \item Créer un catalogue de formations structuré en modules et leçons, selon un modèle freemium.
    \item Intégrer un système de quiz à correction automatique, avec surveillance légère (proctoring).
    \item Développer un tableau de bord de suivi de progression.
    \item Générer des certificats de réussite.
    \item Intégrer un chatbot pédagogique intelligent (ARIA) basé sur un LLM.
    \item Mettre en place un système de paiement pour des formations premium.
    \item Assurer l'internationalisation (Français/Arabe).
    \item Fournir une interface d'administration.
    \item Assurer un niveau de sécurité applicative renforcé (RLS, CSP/HSTS, audit, anti brute-force).
\end{itemize}

\section{Étude de l'existant}
Avant d'entamer la conception, une étude des solutions e-learning existantes a été menée afin d'identifier les bonnes pratiques ainsi que les lacunes à combler. Cette étude a porté sur des plateformes généralistes (Coursera, OpenClassrooms) et sur des solutions plus spécialisées en IA.

Cette analyse a permis de dégager plusieurs constats :
\begin{itemize}
    \item Les plateformes généralistes offrent un catalogue large mais un accompagnement pédagogique peu personnalisé.
    \item Peu de solutions proposent un contenu en arabe avec support RTL natif.
    \item L'accompagnement par un assistant conversationnel intégré au parcours reste rare, en particulier dans un contexte francophone/arabophone.
    \item La dimension sécurité applicative (protection des données, contrôle d'accès) est rarement mise en avant dans la communication de ces plateformes.
\end{itemize}
Ces constats ont directement orienté les choix fonctionnels de \textit{IAAI eLearning 101} : bilinguisme natif, chatbot pédagogique ARIA, et durcissement sécurité dès la conception.

\section{Planification du projet}
Le projet a été réalisé sur une période de deux mois, du 08 mai 2026 au 08 juillet 2026, suivie d'une phase d'amélioration continue.

\begin{figure}[H]
\centering
\begin{ganttchart}[hgrid, vgrid, x unit=0.75cm, y unit chart=0.7cm, title height=1, bar height=0.6, group height=0.7, progress label text={}, bar/.append style={fill=blue!40}, group/.append style={fill=gray!40}]{1}{9}
\gantttitle{Planification du projet IAAI eLearning 101}{9} \\
\gantttitlelist{S1,S2,S3,S4,S5,S6,S7,S8,S9}{1} \\
\ganttgroup{Analyse et conception}{1}{2} \\
\ganttbar{Analyse des besoins}{1}{1} \\
\ganttbar{Etude de l'existant}{1}{1} \\
\ganttbar{Diagrammes MERISE/UML}{2}{2} \\
\ganttbar{Conception BDD}{2}{2} \\
\ganttgroup{Developpement}{3}{7} \\
\ganttbar{Configuration}{3}{3} \\
\ganttbar{Authentification}{3}{4} \\
\ganttbar{Gestion formations}{4}{5} \\
\ganttbar{Quiz, evaluations et proctoring}{5}{6} \\
\ganttbar{Suivi progression}{5}{6} \\
\ganttbar{Certificats}{6}{6} \\
\ganttbar{Paiement Stripe}{6}{7} \\
\ganttgroup{Chatbot ARIA}{6}{7} \\
\ganttbar{Integration Gemini}{6}{7} \\
\ganttbar{Developpement ARIA}{7}{7} \\
\ganttgroup{Securite}{7}{8} \\
\ganttbar{RLS, CSP/HSTS, audit log}{7}{8} \\
\ganttgroup{Tests et deploiement}{8}{9} \\
\ganttbar{Tests fonctionnels et securite}{8}{8} \\
\ganttbar{Deploiement final}{9}{9} \\
\ganttgroup{Documentation}{7}{9} \\
\ganttbar{Redaction rapport}{7}{9} \\
\ganttmilestone{Fin du projet}{9}
\end{ganttchart}
\caption{Diagramme de Gantt du projet IAAI eLearning 101}
\label{fig:gantt}
\end{figure}

\section{Conclusion}
Ce chapitre a permis de situer le cadre du stage en présentant l'organisme d'accueil, l'étude de l'existant, le cahier des charges et la planification retenue.

\chapter{Analyse et conception}
\section{Analyse des besoins}
\subsection{Identification des acteurs}
\begin{itemize}
    \item \textbf{Apprenant :} utilisateur principal de la plateforme.
    \item \textbf{Administrateur :} responsable de la gestion globale.
    \item \textbf{Assistant ARIA :} assistant pédagogique intelligent.
\end{itemize}

\subsection{Besoins fonctionnels}
\begin{table}[H]
\centering
\caption{Résumé des besoins fonctionnels}
\begin{tabularx}{\textwidth}{|c|X|}
\hline
\textbf{Code} & \textbf{Description} \\
\hline
BF1 & Authentification et gestion des utilisateurs \\
\hline
BF2 & Gestion de formation (modèle freemium) \\
\hline
BF3 & Gestion des modules et des leçons \\
\hline
BF4 & Gestion des quiz et évaluations \\
\hline
BF5 & Suivi de progression \\
\hline
BF6 & Génération des certificats \\
\hline
BF7 & Paiement sécurisé via Stripe \\
\hline
BF8 & Chatbot pédagogique intelligent ARIA \\
\hline
BF9 & Gestion de la communauté \\
\hline
BF10 & Notifications système \\
\hline
BF11 & Recommandations pédagogiques \\
\hline
BF12 & Surveillance légère des quiz (proctoring) \\
\hline
\end{tabularx}
\end{table}

\subsection{Besoins non fonctionnels}
\begin{itemize}
    \item \textbf{Sécurité :} protection des données, authentification sécurisée, politiques RLS strictes, verrouillage anti brute-force, en-têtes CSP/HSTS, journal d'audit immuable.
    \item \textbf{Performance :} temps de réponse rapide.
    \item \textbf{Disponibilité :} accès permanent aux services.
    \item \textbf{Compatibilité :} navigateurs modernes.
    \item \textbf{Responsive Design :} adaptation multi-écrans.
    \item \textbf{Maintenabilité :} architecture modulaire.
    \item \textbf{Évolutivité :} ajout de nouvelles fonctionnalités.
    \item \textbf{Internationalisation :} français et arabe (RTL).
    \item \textbf{Fiabilité :} intégrité des données PostgreSQL.
\end{itemize}

\section{Modélisation MERISE et UML}
La conception de la base de données a été menée selon la méthode MERISE (Modèle Conceptuel puis Logique de Données), complétée par une modélisation UML pour la représentation des aspects fonctionnels et dynamiques du système.

\subsection{Diagramme de cas d'utilisation}
\begin{figure}[H]
    \centering
    \includegraphics[height=0.8\textheight]{images/cas-use.png}
    \caption{Diagramme de cas d'utilisation}
    \label{fig:usecase}
\end{figure}

\subsection{Diagramme de classes}
\begin{figure}[H]
    \centering
    \includegraphics[width=\textwidth]{images/classe.png}
    \caption{Diagramme de classe}
    \label{fig:class}
\end{figure}

\section{Diagrammes de séquence}
\subsection{Authentification}
\begin{figure}[H]
    \centering
    \includegraphics[width=0.95\textwidth]{images/séquence1.png}
    \caption{Diagramme de séquence d'authentification}
    \label{fig:seq_auth}
\end{figure}

\subsection{Inscription et paiement}
\begin{figure}[H]
    \centering
    \includegraphics[width=1\textwidth]{images/séquence2.png}
    \caption{Diagramme de séquence d'inscription et paiement}
    \label{fig:seq_payment}
\end{figure}

\subsection{Suivi de progression}
\begin{figure}[H]
    \centering
    \includegraphics[width=0.95\textwidth]{images/séquence3.png}
    \caption{Diagramme de séquence du suivi de progression}
    \label{fig:seq_progress}
\end{figure}

\subsection{Chatbot ARIA et certificat}
\begin{figure}[H]
    \centering
    \includegraphics[width=0.95\textwidth]{images/séquence4.png}
    \caption{Diagramme de séquence du chatbot ARIA}
    \label{fig:seq_aria}
\end{figure}

\subsection{Diagramme de déploiement}
\begin{figure}[H]
    \centering
    \includegraphics[width=0.95\textwidth]{images/déploiement.png}
    \caption{Diagramme de déploiement}
    \label{fig:deployment}
\end{figure}

\section{Conception de la base de données}
\subsection{Modèle Conceptuel de Données (MCD)}
\begin{figure}[H]
    \centering
    \includegraphics[width=\textwidth]{images/mcd.png}
    \caption{Modèle Conceptuel de Données}
    \label{fig:mcd}
\end{figure}

\subsection{Modèle Logique de Données (MLD)}
\begin{figure}[H]
    \centering
    \includegraphics[width=\textwidth]{images/MLD.png}
    \caption{Modèle Logique de Données}
    \label{fig:mld}
\end{figure}

\chapter{Réalisation}
\section{Introduction}
Cette phase présente l'environnement de développement, les choix technologiques et l'implémentation des fonctionnalités.

\section{Environnement de développement}
\subsection{Matériel utilisé}
\begin{itemize}
    \item Processeur : Intel Core i7 (12\ieme génération)
    \item Mémoire vive : 32 Go DDR5
    \item Stockage : SSD NVMe 1 To
    \item Système d'exploitation : Windows 11
\end{itemize}

\subsection{Logiciels utilisés}
\begin{table}[H]
\centering
\caption{Environnement logiciel}
\begin{tabular}{|l|l|p{6cm}|}
\hline
\textbf{Logiciel} & \textbf{Version} & \textbf{Usage} \\
\hline
Visual Studio Code & 1.90 & Développement et débogage \\
\hline
Node.js & 20.12 & Exécution JavaScript \\
\hline
npm & 10.5 & Gestionnaire de paquets \\
\hline
Git & 2.43 & Gestion de versions \\
\hline
Postman & 10.24 & Test des API \\
\hline
Docker Desktop & 4.28 & Conteneurisation \\
\hline
Figma & Web & Maquettage UI/UX \\
\hline
Google Chrome & 125 & Tests fonctionnels \\
\hline
\end{tabular}
\end{table}

\section{Choix technologiques}
\begin{table}[H]
\centering
\caption{Technologies utilisées}
\begin{tabularx}{\textwidth}{|>{\centering\arraybackslash}p{2.7cm}|X|p{3.6cm}|}
\hline
\textbf{Technologie} & \textbf{Rôle} & \textbf{Justification} \\
\hline
React 18 & Frontend & Interface moderne et performante \\
\hline
Vite 5 & Build Tool & Compilation rapide \\
\hline
Tailwind CSS & Design UI & Développement rapide et responsive \\
\hline
Zustand & State management & Gestion légère de l'état d'authentification et de progression \\
\hline
Supabase & Backend & Authentification, base de données et fonctions intégrées \\
\hline
PostgreSQL & Base de données & Robustesse et fiabilité \\
\hline
Row Level Security (RLS) & Sécurité des données & Contrôle d'accès au niveau ligne, natif à PostgreSQL/Supabase \\
\hline
Manim & Animations pédagogiques & Vidéos explicatives de haute qualité \\
\hline
Gemini 2.5 Flash & IA générative & Chatbot pédagogique ARIA, réponses rapides et pertinentes \\
\hline
Stripe & Paiement & Solution sécurisée et fiable \\
\hline
Vercel & Hébergement & Déploiement continu du frontend, en-têtes de sécurité (CSP/HSTS) \\
\hline
i18next & Traduction & Support multilingue et RTL \\
\hline
\end{tabularx}
\end{table}
\section{Conception de l'interface utilisateur}
\subsection{Approche de conception}

La conception de l'interface utilisateur de la plateforme IAAI eLearning 101 s'est appuyée sur une approche centrée sur l'utilisateur, en tenant compte des besoins spécifiques des apprenants en Intelligence Artificielle. L'objectif était de créer une expérience intuitive, engageante et adaptée aux différents niveaux de compétence des utilisateurs.

Le choix du design a été guidé par plusieurs principes fondamentaux :
\begin{itemize}
    \item \textbf{Simplicité et clarté :} Une interface épurée facilitant la navigation et la compréhension des contenus.
    \item \textbf{Accessibilité :} Conforme aux standards d'accessibilité pour un public large.
    \item \textbf{Responsive Design :} Adaptation optimale sur tous les types d'écrans (desktop, tablette, mobile).
    \item \textbf{Cohérence visuelle :} Utilisation d'une palette de couleurs et d'un style uniforme.
    \item \textbf{Feedback utilisateur :} Indicateurs visuels clairs pour les actions et interactions.
\end{itemize}

\subsection{Maquettes de l'interface}

Les maquettes présentées ci-dessous illustrent les principales interfaces de la plateforme, conçues pour offrir une expérience utilisateur fluide et cohérente.

\subsubsection{Landing Page}
\begin{figure}[H]
    \centering
    \includegraphics[width=0.95\textwidth]{images/maquette-landing.png}
    \caption{Maquette de la page d'accueil}
    \label{fig:maquette_landing}
\end{figure}

La landing page constitue le point d'entrée principal de la plateforme. Elle présente de manière claire et attrayante la proposition de valeur de IAAI eLearning 101, les formations disponibles, ainsi que les avantages de l'apprentissage de l'IA. Le design met en avant les éléments d'appel à l'action (CTA) pour encourager l'inscription et l'exploration des contenus.

\subsubsection{Page d'inscription}
\begin{figure}[H]
    \centering
    \includegraphics[width=0.8\textwidth]{images/maquette-register.png}
    \caption{Maquette de la page d'inscription}
    \label{fig:maquette_register}
\end{figure}

La page d'inscription a été conçue pour être simple et rapide, minimisant les frictions lors de la création de compte. Le formulaire est structuré de manière logique avec des indications claires pour chaque champ, et inclut une validation en temps réel pour améliorer l'expérience utilisateur.

\subsubsection{Page de connexion}
\begin{figure}[H]
    \centering
    \includegraphics[width=0.8\textwidth]{images/maquette-login.png}
    \caption{Maquette de la page de connexion}
    \label{fig:maquette_login}
\end{figure}

La page de connexion suit les mêmes principes de simplicité que la page d'inscription, avec un formulaire épuré et des options de récupération de mot de passe facilement accessibles. L'interface inclut également la possibilité de connexion via des réseaux sociaux pour faciliter l'accès.

\subsubsection{Tableau de bord (Dashboard)}
\begin{figure}[H]
    \centering
    \includegraphics[width=0.95\textwidth]{images/maquette-dashboard.png}
    \caption{Maquette du tableau de bord}
    \label{fig:maquette_dashboard}
\end{figure}

Le tableau de bord constitue le centre de contrôle personnel de l'apprenant. Il présente une vue synthétique de la progression, des formations en cours, des certificats obtenus et des statistiques d'apprentissage. Le design utilise des cartes et des graphiques pour rendre les données facilement digestibles et motivantes.

\subsubsection{Parcours IA}
\begin{figure}[H]
    \centering
    \includegraphics[width=0.9\textwidth]{images/maquette-parcours-ai.png}
    \caption{Maquette du parcours d'apprentissage IA}
    \label{fig:maquette_parcours}
\end{figure}

L'interface du parcours d'apprentissage organise les modules et leçons de manière progressive, avec une barre de progression visuelle pour motiver l'apprenant. Chaque leçon est présentée avec ses objectifs, sa durée estimée et son statut de complétion, permettant une navigation intuitive dans le contenu.

\subsubsection{Catalogue des cours}
\begin{figure}[H]
    \centering
    \includegraphics[width=0.95\textwidth]{images/maquette-catalogue.png}
    \caption{Maquette du catalogue des formations}
    \label{fig:maquette_catalogue}
\end{figure}

Le catalogue des cours utilise une présentation en cartes avec des filtres et une recherche intuitive pour aider les apprenants à trouver les formations adaptées à leurs besoins. Chaque carte présente les informations essentielles : titre, description, niveau, durée et indicateur de contenu premium ou gratuit.

\subsubsection{Profil utilisateur}
\begin{figure}[H]
    \centering
    \includegraphics[width=0.8\textwidth]{images/maquette-profile.png}
    \caption{Maquette du profil utilisateur}
    \label{fig:maquette_profile}
\end{figure}

La page de profil permet à l'utilisateur de gérer ses informations personnelles, ses préférences et son historique d'apprentissage. L'interface est organisée en sections claires avec des formulaires de modification simples et sécurisés.

\subsubsection{Communauté}
\begin{figure}[H]
    \centering
    \includegraphics[width=0.95\textwidth]{images/maquette-communaute.png}
    \caption{Maquette de l'espace communautaire}
    \label{fig:maquette_communaute}
\end{figure}

L'espace communautaire favorise les échanges entre apprenants à travers une interface de type réseau social. Les utilisateurs peuvent publier des questions, partager leurs expériences et interagir avec les autres membres de la communauté, créant ainsi un environnement d'apprentissage collaboratif.

\subsubsection{Certificat}
\begin{figure}[H]
    \centering
    \includegraphics[width=0.8\textwidth]{images/maquette-certificat.png}
    \caption{Maquette du certificat de réussite}
    \label{fig:maquette_certificat}
\end{figure}

L'interface de certificat présente les certificats obtenus par l'apprenant avec un design professionnel et sécurisé. Chaque certificat affiche les détails de la formation validée, la date d'obtention et un code de vérification unique pour garantir l'authenticité.

\subsubsection{Notifications}
\begin{figure}[H]
    \centering
    \includegraphics[width=0.8\textwidth]{images/maquette-notifications.png}
    \caption{Maquette du centre de notifications}
    \label{fig:maquette_notifications}
\end{figure}

Le centre de notifications regroupe toutes les alertes et informations importantes pour l'apprenant : rappels de formations, mises à jour de contenu, résultats d'évaluations et messages de la communauté. L'interface permet de filtrer et de gérer les notifications selon leurs types et leur importance.

\subsubsection{Paramètres}
\begin{figure}[H]
    \centering
    \includegraphics[width=0.8\textwidth]{images/maquette-parametres.png}
    \caption{Maquette de la page des paramètres}
    \label{fig:maquette_parametres}
\end{figure}

La page des paramètres offre un contrôle complet sur les préférences de l'utilisateur : langue (français/arabe), notifications, confidentialité, et options d'accessibilité. L'interface est organisée de manière hiérarchique pour faciliter la navigation entre les différentes catégories de paramètres.

\subsubsection{Page d'erreur}
\begin{figure}[H]
    \centering
    \includegraphics[width=0.8\textwidth]{images/maquette-error.png}
    \caption{Maquette de la page d'erreur}
    \label{fig:maquette_error}
\end{figure}

La page d'erreur a été conçue pour être informative et helpful, même en cas de problème technique. Elle présente un message clair expliquant l'erreur, des suggestions de résolution et des boutons de navigation alternatifs pour aider l'utilisateur à continuer son expérience.

\subsection{Choix du design et justification}

Le design global de la plateforme repose sur une approche minimaliste et moderne, utilisant une palette de couleurs cohérente avec l'identité visuelle de l'IAAI Academy. Le choix d'une interface épurée vise à réduire la charge cognitive de l'apprenant, lui permettant de se concentrer sur le contenu pédagogique plutôt que sur la navigation complexe.

L'utilisation de composants UI réutilisables et d'une grille responsive garantit une expérience cohérente sur tous les appareils, tandis que les micro-interactions et les animations subtiles améliorent l'engagement sans distraire de l'apprentissage. Le support du mode RTL pour l'arabe et l'internationalisation complète démontrent l'attention portée à l'accessibilité et à l'inclusivité de la plateforme.

\section{Sécurité applicative}
Au-delà des fonctionnalités pédagogiques, une part importante du travail a porté sur le durcissement de la sécurité de la plateforme, en particulier suite à un audit fonctionnel ayant révélé plusieurs failles.

\subsection{Politiques Row Level Security (RLS)}
L'ensemble des tables sensibles (profils, formations, tentatives de quiz, certificats) est protégé par des politiques RLS strictes sur PostgreSQL, restreignant l'écriture aux seuls rôles autorisés (par exemple, les insertions dans la table \texttt{certificates} sont désormais interdites aux rôles \texttt{anon}/\texttt{authenticated} et ne passent que par des fonctions RPC dédiées).

\subsection{Fonctions RPC sécurisées (SECURITY DEFINER)}
Certaines opérations sensibles ne sont plus exécutées directement depuis le client, mais via des fonctions PostgreSQL en mode \texttt{SECURITY DEFINER} :
\begin{itemize}
    \item \texttt{issue\_certificate()} : génération et validation de l'éligibilité au certificat.
    \item \texttt{verify\_certificate()} : exposition contrôlée des seuls champs publics d'un certificat, pour la vérification externe.
    \item \texttt{is\_login\_locked()} / \texttt{record\_login\_attempt()} : gestion du verrouillage de connexion.
\end{itemize}

\subsection{Protection contre les attaques par force brute}
Un mécanisme de verrouillage de connexion a été mis en place à partir d'une table \texttt{login\_attempts}, sans aucune politique RLS ouverte, uniquement accessible via les fonctions \texttt{SECURITY DEFINER} citées ci-dessus.

\subsection{En-têtes de sécurité et audit}
\begin{itemize}
    \item Mise en place d'en-têtes CSP (Content Security Policy) et HSTS stricts au niveau de la configuration Vercel (\texttt{vercel.json}).
    \item Journal d'audit immuable (\texttt{admin\_audit\_log}) : aucune politique \texttt{UPDATE}/\texttt{DELETE} n'est définie, garantissant la traçabilité des actions d'administration (suppression et invitation d'utilisateurs notamment).
    \item Suppression du jeton d'accès (\texttt{accessToken}) du stockage local (Zustand persist), pour réduire la surface d'exposition en cas de faille XSS.
\end{itemize}

\subsection{Travaux de sécurité restant à finaliser}
Certains points nécessitent un accès direct au tableau de bord Supabase et n'ont pas encore été traités au moment de la rédaction de ce rapport :
\begin{itemize}
    \item Vérification par export (\texttt{db dump}) des politiques RLS sur les tables non versionnées.
    \item Activation de la protection contre les mots de passe compromis (HaveIBeenPwned) dans Supabase Auth.
    \item Vérification de l'activation du \textit{secret scanning} GitHub.
\end{itemize}

\section{Chatbot pédagogique ARIA}
\subsection{Principe de fonctionnement}
ARIA est un chatbot pédagogique conçu pour accompagner l'apprenant tout au long de son parcours sur la plateforme. Contrairement à une architecture de type RAG (Retrieval-Augmented Generation), ARIA repose sur un appel direct à l'API Gemini, orchestré par un prompt système décrivant son rôle, le contexte de la leçon en cours et les données de progression de l'apprenant.

\begin{figure}[H]
    \centering
    \begin{tikzpicture}[
        node distance=1cm and 2cm,
        box/.style={rectangle, rounded corners, draw=black, thick, fill=blue!10, minimum width=4.5cm, minimum height=1.1cm, align=center, font=\small},
        boxalt/.style={rectangle, rounded corners, draw=black, thick, fill=orange!15, minimum width=4.5cm, minimum height=1.1cm, align=center, font=\small},
        boxedge/.style={rectangle, rounded corners, draw=black, thick, fill=green!10, minimum width=4.5cm, minimum height=1.1cm, align=center, font=\small},
        arrow/.style={-{Latex[length=2.5mm]}, thick}
    ]
    \node[box] (user) {Apprenant\\\textit{(pose une question)}};
    \node[boxedge, below=of user] (edge) {Interface Web\\\textit{(envoi de la question)}};
    \node[box, below=of edge] (context) {Récupération du contexte\\(leçon en cours, progression)};
    \node[boxalt, below=of context] (prompt) {Construction du prompt système};
    \node[box, below=of prompt] (llm) {Appel à l'API Gemini 2.5 Flash};
    \node[boxedge, below=of llm] (response) {Réponse générée\\renvoyée à l'apprenant};
    \draw[arrow] (user) -- (edge);
    \draw[arrow] (edge) -- (context);
    \draw[arrow] (context) -- (prompt);
    \draw[arrow] (prompt) -- (llm);
    \draw[arrow] (llm) -- (response);
    \draw[arrow] (response.west) -- ++(-2.7,0) |- (user.west);
    \end{tikzpicture}
    \caption{Architecture du chatbot ARIA}
    \label{fig:aria-architecture}
\end{figure}

\subsection{Fonctionnalités couvertes}
\begin{itemize}
    \item Répondre aux questions de l'apprenant sur les notions abordées dans la leçon en cours.
    \item Fournir un accompagnement contextualisé à la progression de l'apprenant.
    \item Assister l'apprenant lors de la consultation et du partage de son certificat.
\end{itemize}

\subsection{Limites actuelles}
N'étant pas adossé à une base de connaissances vectorisée, ARIA ne recherche pas d'information dans le contenu détaillé des leçons au-delà de ce qui est transmis explicitement dans le prompt système. Cette limite constitue une perspective d'évolution identifiée pour une version future de la plateforme.

\section{Production des vidéos pédagogiques}
Afin de faciliter la compréhension de notions abstraites d'Intelligence Artificielle, des vidéos d'animation ont été produites pour les leçons du Module 1, selon un pipeline dédié.

\subsection{Charte graphique}
Les animations suivent la charte graphique « IAAI Nebula », avec deux couleurs principales : violet (\texttt{PURPLE\_MAIN \#7C3AED}) et rose (\texttt{PINK\_ACCENT \#EC4899}).

\subsection{Pipeline de production}
\begin{enumerate}
    \item \textbf{Diagnostic du contenu} de la leçon à illustrer.
    \item \textbf{Génération de la narration} en français via \texttt{edge-tts} (voix \texttt{fr-FR-DeniseNeural}).
    \item \textbf{Rendu des animations} avec la bibliothèque Manim.
    \item \textbf{Fusion audio/vidéo} avec \texttt{ffmpeg}.
    \item \textbf{Export} des vidéos finales dans \texttt{public/videos/}.
\end{enumerate}

\subsection{Lecture des vidéos}
Le composant \texttt{VideoPlayer.jsx} prend en charge la lecture des vidéos (intégration YouTube ou lecture HTML5 native), et déclenche automatiquement la validation de la leçon (\texttt{markLessonComplete}) à la fin de la lecture.

\section{Réalisation des fonctionnalités}
\subsection{Page d'accueil}
\begin{figure}[H]
    \centering
    \includegraphics[height=0.9\textheight,keepaspectratio]{images/landing_page.png}
    \caption{Page d'accueil}
    \label{fig:landing_page}
\end{figure}

\subsection{Authentification}
\begin{figure}[H]
\centering
\includegraphics[width=0.8\textwidth]{images/register.png}
\caption{Page d'inscription}
    \label{fig:register}
\end{figure}

\begin{figure}[H]
\centering
\includegraphics[width=0.8\textwidth]{images/login.png}
\caption{Page de connexion}
    \label{fig:login}
\end{figure}

\subsection{Catalogue des formations}
\begin{figure}[H]
\centering
\includegraphics[width=0.95\textwidth]{images/catalogue.png}
\caption{Catalogue des formations}
    \label{fig:catalogue}
\end{figure}

\subsection{Parcours pédagogique}
\begin{figure}[H]
\centering
\includegraphics[width=0.9\textwidth]{images/course-view.png}
\caption{Parcours pédagogique}
    \label{fig:course}
\end{figure}

\subsection{Quiz et évaluations}
\begin{figure}[H]
\centering
\includegraphics[width=0.85\textwidth]{images/page-quiz.jpg}
\caption{Interface de quiz (avec surveillance légère / proctoring)}
\label{fig:quiz}
\end{figure}

\subsection{Tableau de bord}
\begin{figure}[H]
\centering
\includegraphics[width=0.95\textwidth]{images/dashboard.png}
\caption{Tableau de bord apprenant}
    \label{fig:dashboard}
\end{figure}

\subsection{Certificats}
\begin{figure}[H]
\centering
\includegraphics[width=0.8\textwidth]{images/certificate.png}
\caption{Certificat généré}
    \label{fig:certificate}
\end{figure}

\subsection{Paiement Stripe}
\begin{figure}[H]
\centering
\includegraphics[width=0.8\textwidth]{images/stripe-payment.png}
\caption{Paiement sécurisé avec Stripe}
    \label{fig:stripe}
\end{figure}

\subsection{Administration}
\begin{figure}[H]
\centering
\includegraphics[width=0.95\textwidth]{images/admin-dashboard.png}
\caption{Interface d'administration}
    \label{fig:admin}
\end{figure}

\subsection{Internationalisation}
\begin{figure}[H]
\centering
\includegraphics[width=0.5\textwidth]{images/arabic-interface.png}
\caption{Version arabe de la plateforme}
\label{fig:arabic}
\end{figure}

\section{Déploiement}
\subsection{Frontend Vercel}
\begin{figure}[H]
\centering
\includegraphics[width=0.85\textwidth]{images/vercel_dashboard.png}
\caption{Tableau de bord Vercel}
\label{fig:vercel_dashboard}
\end{figure}

\subsection{Backend Supabase}
\begin{figure}[H]
\centering
\includegraphics[width=0.85\textwidth]{images/supabase_dashboard.png}
\caption{Tableau de bord Supabase}
\label{fig:supabase_dashboard}
\end{figure}

\subsection{Base PostgreSQL}
\begin{figure}[H]
\centering
\includegraphics[width=0.85\textwidth]{images/postgresql_database.png}
\caption{Base de données PostgreSQL}
\label{fig:postgres_database}
\end{figure}

\subsection{API Gemini}
\begin{figure}[H]
\centering
\includegraphics[width=0.85\textwidth]{images/gemini_api.png}
\caption{Configuration API Gemini}
\label{fig:gemini_api}
\end{figure}

\subsection{Stripe}
\begin{figure}[H]
\centering
\includegraphics[width=0.85\textwidth]{images/stripe_dashboard.png}
\caption{Interface Stripe}
\label{fig:stripe_dashboard}
\end{figure}

\chapter{Tests et Validation}
\section{Introduction}
Plusieurs campagnes de tests ont été effectuées pour vérifier la conformité de la plateforme aux spécifications, ainsi que la robustesse des mesures de sécurité mises en place.

\section{Tests fonctionnels}
\subsection{Authentification}
\begin{table}[H]
\centering
\caption{Tests d'authentification}
\begin{tabularx}{\textwidth}{|p{4cm}|X|c|}
\hline
\textbf{Cas de test} & \textbf{Résultat attendu} & \textbf{Statut} \\
\hline
Création compte & Création + email confirmation & Succès \\
\hline
Connexion valide & Accès tableau de bord & Succès \\
\hline
Mot de passe incorrect & Message d'erreur & Succès \\
\hline
Réinitialisation & Lien réinitialisation & Succès \\
\hline
Tentatives multiples échouées & Verrouillage temporaire du compte & Succès \\
\hline
\end{tabularx}
\end{table}

\subsection{Catalogue}
\begin{table}[H]
\centering
\caption{Tests catalogue}
\begin{tabularx}{\textwidth}{|p{4cm}|X|c|}
\hline
\textbf{Cas de test} & \textbf{Résultat attendu} & \textbf{Statut} \\
\hline
Affichage catalogue & Chargement formations & Succès \\
\hline
Recherche mot-clé & Résultats pertinents & Succès \\
\hline
Filtrage catégorie & Filtrage correct & Succès \\
\hline
Accès formation & Détail du cours & Succès \\
\hline
\end{tabularx}
\end{table}

\subsection{Quiz et proctoring}
\begin{table}[H]
\centering
\caption{Tests quiz et surveillance}
\begin{tabularx}{\textwidth}{|p{4cm}|X|c|}
\hline
\textbf{Cas de test} & \textbf{Résultat attendu} & \textbf{Statut} \\
\hline
Soumission quiz & Réponses enregistrées & Succès \\
\hline
Calcul score & Score correct & Succès \\
\hline
Seuil réussite & Statut réussi/échoué & Succès \\
\hline
Historique & Sauvegarde correcte & Succès \\
\hline
Sortie d'onglet pendant le quiz & Quiz invalidé immédiatement & Succès \\
\hline
Violation visible côté admin & Journal par étudiant et par quiz & Succès \\
\hline
\end{tabularx}
\end{table}

\subsection{Certificats}
\begin{table}[H]
\centering
\caption{Tests certificats}
\begin{tabularx}{\textwidth}{|p{4cm}|X|c|}
\hline
\textbf{Cas de test} & \textbf{Résultat attendu} & \textbf{Statut} \\
\hline
Validation conditions & Vérification correcte & Succès \\
\hline
Génération PDF & Création certificat & Succès \\
\hline
Téléchargement & Fichier accessible & Succès \\
\hline
Code vérification & Code unique & Succès \\
\hline
Vérification publique (\texttt{/verify/:numero}) & Page accessible, champs publics uniquement & Succès \\
\hline
\end{tabularx}
\end{table}

\subsection{Paiement Stripe}
\begin{table}[H]
\centering
\caption{Tests paiement}
\begin{tabularx}{\textwidth}{|p{4cm}|X|c|}
\hline
\textbf{Cas de test} & \textbf{Résultat attendu} & \textbf{Statut} \\
\hline
Session Stripe & Session créée & Succès \\
\hline
Paiement accepté & Déblocage premium & Succès \\
\hline
Paiement refusé & Message d'erreur & Succès \\
\hline
Webhook & Mise à jour statut & Succès \\
\hline
\end{tabularx}
\end{table}

\subsection{Chatbot ARIA}
\begin{table}[H]
\centering
\caption{Tests ARIA}
\begin{tabularx}{\textwidth}{|p{5cm}|X|c|}
\hline
\textbf{Question} & \textbf{Résultat observé} & \textbf{Statut} \\
\hline
Qu'est-ce que le ML ? & Réponse cohérente & Succès \\
\hline
Explique le Deep Learning & Réponse contextualisée & Succès \\
\hline
IA faible / IA forte & Réponse correcte & Succès \\
\hline
\end{tabularx}
\end{table}

\section{Tests de sécurité}
\begin{table}[H]
\centering
\caption{Tests de sécurité}
\begin{tabularx}{\textwidth}{|p{5cm}|X|c|}
\hline
\textbf{Cas de test} & \textbf{Résultat attendu} & \textbf{Statut} \\
\hline
Insertion directe dans \texttt{certificates} par un client anon/authenticated & Rejetée par la politique RLS & Succès \\
\hline
Lecture/écriture sur \texttt{login\_attempts} depuis le client & Rejetée (aucune politique RLS ouverte) & Succès \\
\hline
Modification/suppression d'une entrée \texttt{admin\_audit\_log} & Rejetée (pas de politique UPDATE/DELETE) & Succès \\
\hline
Tentatives de connexion répétées & Verrouillage déclenché après le seuil configuré & Succès \\
\hline
En-têtes de réponse HTTP & Présence de CSP et HSTS & Succès \\
\hline
Accès à \texttt{/verify/:numero} avec un numéro invalide & Erreur contrôlée, pas de fuite de données & Succès \\
\hline
\end{tabularx}
\end{table}

\section{Validation du déploiement}
\begin{table}[H]
\centering
\caption{Validation composants}
\begin{tabular}{|l|c|}
\hline
Composant & Statut \\
\hline
Frontend Vercel & Opérationnel \\
\hline
Backend Supabase & Opérationnel \\
\hline
Base PostgreSQL & Opérationnelle \\
\hline
Gemini API & Opérationnelle \\
\hline
Stripe API & Opérationnelle \\
\hline
\end{tabular}
\end{table}

\chapter*{Conclusion Générale}
\addcontentsline{toc}{chapter}{Conclusion Générale}

Ce stage de fin d'études réalisé au sein de l'IAAI Academy a permis la conception et le développement de la plateforme \textit{IAAI eLearning 101}, une solution d'apprentissage en ligne dédiée à la formation dans les domaines de l'intelligence artificielle.

L'objectif principal était de mettre en place une plateforme moderne permettant aux apprenants d'accéder à des contenus pédagogiques structurés, de suivre leur progression, de réaliser des évaluations et d'obtenir des certifications numériques. Un autre objectif majeur consistait à intégrer un assistant pédagogique intelligent capable d'accompagner les utilisateurs tout au long de leur parcours d'apprentissage.

Une architecture moderne reposant sur React, Supabase et PostgreSQL a été mise en œuvre. Les différentes fonctionnalités prévues ont été développées avec succès : authentification sécurisée, gestion des formations en modèle freemium, suivi de progression, quiz avec surveillance légère (proctoring), certificats, paiements Stripe, internationalisation et chatbot ARIA basé sur l'API Gemini. Un travail conséquent de durcissement sécurité (politiques RLS strictes, protection anti brute-force, en-têtes CSP/HSTS, journal d'audit immuable) a également été mené suite à un audit fonctionnel, renforçant la fiabilité globale de la plateforme.

Les phases de tests et de validation ont permis de vérifier le bon fonctionnement de l'ensemble des fonctionnalités développées ainsi que la robustesse des mesures de sécurité mises en place. Les résultats obtenus démontrent la robustesse de la solution et sa capacité à offrir une expérience utilisateur fluide et sécurisée.

Au-delà des compétences techniques acquises, ce stage a constitué une expérience professionnelle enrichissante, permettant de renforcer les compétences en gestion de projet, analyse des besoins, conception logicielle et résolution de problèmes complexes.

Plusieurs perspectives d'évolution peuvent être envisagées :
\begin{itemize}
    \item Développement d'une application mobile native.
    \item Intégration de mécanismes de gamification avancés.
    \item Ajout de recommandations pédagogiques basées sur le ML.
    \item Évolution du chatbot ARIA vers une architecture RAG s'appuyant sur une base de connaissances vectorisée.
    \item Intégration de classes virtuelles et fonctionnalités collaboratives.
    \item Mise en place d'outils analytiques avancés.
    \item Finalisation des points de sécurité restants (protection mots de passe compromis, audit des politiques RLS sur les tables non versionnées).
\end{itemize}

En conclusion, ce projet a permis de concrétiser une plateforme e-learning sécurisée et intelligente répondant aux besoins actuels de la formation numérique, constituant une base solide pour le développement futur de solutions pédagogiques innovantes.

\chapter*{Bibliographie}
\addcontentsline{toc}{chapter}{Bibliographie}

\begin{enumerate}
\item React Team, \textit{React Documentation}, Meta Platforms Inc., 2025.
\item Supabase Inc., \textit{Supabase Documentation}, 2025.
\item PostgreSQL Global Development Group, \textit{PostgreSQL Documentation}, Version 15, 2025.
\item Stripe Inc., \textit{Stripe API Documentation}, 2025.
\item Google DeepMind, \textit{Gemini API Documentation}, 2025.
\item Vercel Inc., \textit{Vercel Documentation}, 2025.
\item Tailwind Labs, \textit{Tailwind CSS Documentation}, Version 3, 2025.
\item i18next Contributors, \textit{i18next Documentation}, 2025.
\item Martin Kleppmann, \textit{Designing Data-Intensive Applications}, O'Reilly Media, 2017.
\item Ian Goodfellow, Yoshua Bengio, Aaron Courville, \textit{Deep Learning}, MIT Press, 2016.
\item Stuart Russell et Peter Norvig, \textit{Artificial Intelligence: A Modern Approach}, Pearson, 4ème édition, 2021.
\item OWASP Foundation, \textit{OWASP Top Ten}, 2025.
\end{enumerate}

\chapter*{Webographie}
\addcontentsline{toc}{chapter}{Webographie}

\begin{enumerate}
\item https://react.dev
\item https://supabase.com/docs
\item https://www.postgresql.org/docs
\item https://stripe.com/docs
\item https://ai.google.dev
\item https://vercel.com/docs
\item https://tailwindcss.com/docs
\item https://www.i18next.com
\item https://owasp.org
\end{enumerate}

\end{document}
